% ============================================================
%  GUÍA COMPLEMENTARIA - CÍRCULO UNITARIO E IDENTIDADES
%  Cálculo I - Usach Premium
%
%  Calibración:
%  - Plantilla institucional definida en template/ESTILOS.md.
%  - Progresión desde valores notables hasta identidades,
%    ecuaciones y aplicaciones.
% ============================================================

\documentclass[letterpaper,12pt]{article}

% ============================================================
% PAQUETES
% ============================================================
\usepackage[utf8]{inputenc}
\usepackage[T1]{fontenc}
\usepackage[spanish]{babel}
\usepackage{amsmath,amssymb,amsfonts}
\usepackage{geometry}
\usepackage{fancyhdr}
\usepackage{enumitem}
\usepackage{graphicx}
\usepackage{hyperref}
\usepackage{multicol}
\usepackage{parskip}
\usepackage{xcolor}
\usepackage{tcolorbox}
\usepackage{eso-pic}
\makeatletter
\providecommand{\IfPDFManagementActiveF}[1]{#1}
\makeatother
\usepackage{transparent}
\usepackage{tikz}
\usetikzlibrary{arrows.meta,babel,calc,patterns,shadows.blur}
\usepackage{array}
\usepackage{booktabs}
\usepackage{colortbl}
\usepackage{needspace}
\tcbuselibrary{skins,breakable}

% ============================================================
% GEOMETRÍA Y COLORES
% ============================================================
\geometry{letterpaper,margin=1.5cm,top=3cm,bottom=2.5cm,headheight=2cm}

\definecolor{celesteSuave}{HTML}{F0F7FF}
\definecolor{azulFuerte}{HTML}{1A5276}
\definecolor{azulMedio}{HTML}{2980B9}
\definecolor{turquesaUsach}{HTML}{33CCCC}
\definecolor{enlace}{HTML}{1A5276}
\definecolor{grisTexto}{HTML}{4D5656}
\definecolor{verdeClave}{HTML}{1E8449}

% ============================================================
% INFORMACIÓN DE LA GUÍA
% ============================================================
\newcommand{\logoup}{../../../../../template/images/logo_up.png}
\newcommand{\logousach}{../../../../../template/images/logo_usach.png}
\newcommand{\logofondo}{../../../../../template/images/logo_up.png}

\newcommand{\institucion}{Usach Premium}
\newcommand{\curso}{Cálculo I}
\newcommand{\guiasnumero}{Guía complementaria - Círculo unitario}
\newcommand{\semestre}{1er. Semestre de 2026}
\newcommand{\preparadopor}{Usach Premium.}
\newcommand{\webinstitucion}{www.usachpremium.cl}
\newcommand{\instagraminstitucion}{https://www.instagram.com/usach.premium}
\newcommand{\transparencia}{0.045}
\newcommand{\fraseportada}{"Por y para estudiantes"}

\newcommand{\seno}{\operatorname{sen}}
\newcommand{\R}{\mathbb{R}}

\newcommand{\listacontenidos}{%
    \item[\color{azulFuerte}$\Rightarrow$] Círculo unitario y valores exactos
    \item[\color{azulFuerte}$\Rightarrow$] Identidades pitagóricas
    \item[\color{azulFuerte}$\Rightarrow$] Ecuaciones trigonométricas base
    \item[\color{azulFuerte}$\Rightarrow$] Ejercicios progresivos y desafíos
}

% ============================================================
% HIPERVÍNCULOS, ENCABEZADO Y PIE
% ============================================================
\hypersetup{
    colorlinks=true,
    linkcolor=enlace,
    urlcolor=enlace,
    citecolor=enlace,
    pdfborder={0 0 0},
    pdftitle={Guía complementaria - Círculo unitario e identidades},
    pdfauthor={Usach Premium}
}

\pagestyle{fancy}
\fancyhf{}
\fancyhead[L]{\includegraphics[height=1.2cm]{\logoup}}
\fancyhead[R]{%
    \begin{minipage}[b]{0.47\textwidth}
        \raggedleft
        \fontsize{9pt}{11pt}\selectfont
        \textbf{\color{azulFuerte}\institucion}\\
        \curso\ - Círculo unitario\\
        \textit{\color{gray}@usach.premium}
    \end{minipage}\hspace{0.1cm}%
    \includegraphics[height=1.2cm]{\logousach}%
}
\fancyfoot[L]{\fontsize{9pt}{11pt}\selectfont\textit{\fraseportada}}
\fancyfoot[R]{\fontsize{9pt}{11pt}\selectfont Página \thepage}
\renewcommand{\headrulewidth}{0.4pt}
\renewcommand{\footrulewidth}{0.4pt}

\fancypagestyle{plain}{%
    \fancyhf{}
    \renewcommand{\headrulewidth}{0pt}
    \renewcommand{\footrulewidth}{0pt}
}

% ============================================================
% MARCA DE AGUA Y CAJAS
% ============================================================
\newcommand{\MarcaAgua}{%
    \AddToShipoutPicture{%
        \AtPageCenter{%
            \makebox(0,0){%
                \transparent{\transparencia}%
                \includegraphics[width=0.7\textwidth]{\logofondo}%
            }%
        }%
    }%
}

\newtcolorbox{teoriabox}[1]{%
    colback=celesteSuave,
    colframe=azulFuerte,
    colbacktitle=azulFuerte,
    coltitle=white,
    title=\textbf{#1},
    fonttitle=\bfseries,
    arc=7pt,
    boxrule=1pt,
    enhanced,
    breakable
}

\newtcolorbox{recordatoriobox}[1]{%
    colback=white,
    colframe=azulMedio,
    title=\textbf{#1},
    fonttitle=\bfseries,
    arc=5pt,
    boxrule=0.9pt
}

\newtcolorbox{desafiointegradorbox}[1]{%
    colback=white,
    colframe=azulFuerte,
    colbacktitle=azulFuerte,
    coltitle=white,
    title=\textbf{#1},
    fonttitle=\bfseries,
    arc=8pt,
    boxrule=1.2pt,
    enhanced,
    drop shadow={black!7},
    breakable
}

\newtcolorbox{solbox}[1]{%
    colback=celesteSuave,
    colframe=azulFuerte,
    title=\textbf{#1},
    fonttitle=\bfseries,
    arc=5pt,
    boxrule=0.9pt,
    breakable
}

\newcommand{\tituloSeccion}[1]{%
    \needspace{4\baselineskip}
    \subsection*{\color{azulFuerte}#1}
}

\setlist[enumerate]{itemsep=0.35em,topsep=0.35em}

% ============================================================
\begin{document}
\hypersetup{pageanchor=false}
\thispagestyle{plain}

% ============================================================
% PORTADA
% ============================================================
\AddToShipoutPicture*{%
    \put(54,700){\includegraphics[width=2cm]{\logoup}}
    \put(420,706){\includegraphics[width=5cm]{\logousach}}
}

\vspace*{0.5cm}

\begin{center}
    \color{azulFuerte}
    {\huge\textbf{GUÍA DE EJERCICIOS}}\\[0.5cm]
    {\Huge\textbf{\curso}}\\[0.3cm]
    {\Large\guiasnumero}\\[1.5cm]

    \begin{tcolorbox}[
        colback=celesteSuave,
        colframe=azulFuerte,
        arc=10pt,
        width=0.85\textwidth,
        center,
        boxrule=1.5pt]
        \centering
        \vspace{0.2cm}
        {\Large\textbf{Contenidos}}
        \vspace{0.3cm}
        \color{black}
        \begin{itemize}[leftmargin=2.6cm,itemsep=0.35cm]
            \listacontenidos
        \end{itemize}
        \vspace{0.2cm}
    \end{tcolorbox}
\end{center}

\vspace{1.5cm}

\begin{center}
    {\Large\textbf{\semestre}}\\[0.8cm]
    {\large\textbf{\preparadopor}}\\[1.2cm]
    \textit{\color{grisTexto}La idea no es memorizar una vuelta: es aprender a leerla.}
\end{center}

\vfill

\begin{center}
    {\color{azulFuerte}\hrule height 0.5pt}
    \vspace{0.3cm}
    \begin{minipage}{0.45\textwidth}
        \centering
        \textbf{Instagram:}\\
        \small\href{\instagraminstitucion}{@usach.premium}
    \end{minipage}
    \begin{minipage}{0.45\textwidth}
        \centering
        \textbf{Web:}\\
        \small\href{http://\webinstitucion}{\webinstitucion}
    \end{minipage}
\end{center}

\clearpage
\pagenumbering{arabic}
\hypersetup{pageanchor=true}
\MarcaAgua

% ============================================================
% TEORÍA BREVE
% ============================================================
\section*{\color{azulFuerte}El círculo unitario, en breve}

\begin{teoriabox}{LA IDEA, SIN PALABRAS COMPLICADAS}
\begin{minipage}[c]{0.56\textwidth}
\raggedright
\begin{enumerate}[leftmargin=*,label=\color{azulFuerte}\bfseries\arabic*.,itemsep=0.45em]
    \item Imagina una rueda de radio \(1\), con su centro en el punto \((0,0)\).
    \item Empezamos en el punto de la derecha, \((1,0)\), y giramos el ángulo \(\theta\).
    \item Al detenernos llegamos a un punto:
    \[
        P(\theta)=(\cos\theta,\seno\theta).
    \]
    El coseno indica si estamos a la derecha o a la izquierda; el seno, si estamos arriba o abajo.
    \item Los mismos números se repiten alrededor del círculo. Según dónde quede el punto, pueden ser positivos o negativos.
    \item Una vuelta completa son \(360^\circ\). En radianes, esa misma vuelta se escribe \(2\pi\).
\end{enumerate}
\end{minipage}\hfill
\begin{minipage}[c]{0.39\textwidth}
\centering
\begin{tikzpicture}[scale=2.25,>=Latex]
    \draw[->,gray!70] (-1.25,0)--(1.28,0) node[right] {\scriptsize \(x\)};
    \draw[->,gray!70] (0,-1.25)--(0,1.28) node[above] {\scriptsize \(y\)};
    \draw[azulFuerte,line width=1.1pt] (0,0) circle (1);
    \coordinate (P) at (40:1);
    \draw[azulMedio,line width=1.1pt,->] (0,0)--(P);
    \draw[dashed,gray] (P)--({cos(40)},0);
    \draw[dashed,gray] (P)--(0,{sin(40)});
    \fill[azulFuerte] (P) circle (0.025);
    \draw[azulMedio,->] (0.3,0) arc[start angle=0,end angle=40,radius=0.3];
    \node[azulFuerte] at (0.37,0.12) {\scriptsize \(\theta\)};
    \node[above right,align=left] at (P) {\scriptsize \(P(\theta)\)\\[-1mm]\scriptsize \((\cos\theta,\seno\theta)\)};
    \node[below] at ({cos(40)},0) {\scriptsize \(\cos\theta\)};
    \node[left] at (0,{sin(40)}) {\scriptsize \(\seno\theta\)};
    \node[gray] at (0.55,0.72) {\tiny I};
    \node[gray] at (-0.55,0.72) {\tiny II};
    \node[gray] at (-0.55,-0.72) {\tiny III};
    \node[gray] at (0.55,-0.72) {\tiny IV};
\end{tikzpicture}
\end{minipage}
\end{teoriabox}

\begin{recordatoriobox}{CÓMO LEER CUALQUIER ÁNGULO}
\begin{enumerate}[label=\color{azulFuerte}\bfseries\arabic*.,leftmargin=*,itemsep=0.25em]
    \item Si el ángulo da más de una vuelta, quita vueltas completas hasta quedarte con una sola.
    \item Mira el ángulo pequeño que queda entre la línea y la horizontal. Ese es el \textbf{ángulo de referencia}.
    \item Mira dónde quedó el punto: derecha o izquierda para el coseno; arriba o abajo para el seno.
\end{enumerate}
\end{recordatoriobox}

\begin{teoriabox}{LOS CINCO ÁNGULOS QUE CONVIENE CONOCER}
\centering
\renewcommand{\arraystretch}{1.45}
\small
\begin{tabular}{>{\columncolor{azulFuerte}\color{white}\bfseries}c|ccccc}
\rowcolor{azulFuerte}
\color{white}\(\theta\) &
\color{white}\(0\) &
\color{white}\(\frac{\pi}{6}\) &
\color{white}\(\frac{\pi}{4}\) &
\color{white}\(\frac{\pi}{3}\) &
\color{white}\(\frac{\pi}{2}\)\\
\(\seno\theta\) & \(0\) & \(\frac12\) & \(\frac{\sqrt2}{2}\) & \(\frac{\sqrt3}{2}\) & \(1\)\\
\(\cos\theta\) & \(1\) & \(\frac{\sqrt3}{2}\) & \(\frac{\sqrt2}{2}\) & \(\frac12\) & \(0\)\\
\(\tan\theta\) & \(0\) & \(\frac{\sqrt3}{3}\) & \(1\) & \(\sqrt3\) & no existe\\
\end{tabular}
\end{teoriabox}

\begin{minipage}[t]{0.47\textwidth}
\begin{recordatoriobox}{¿POSITIVO O NEGATIVO?}
\centering
\renewcommand{\arraystretch}{1.25}
\begin{tabular}{c|ccc}
\textbf{Cuadrante} & \(\seno\) & \(\cos\) & \(\tan\)\\ \hline
I   & \(+\) & \(+\) & \(+\)\\
II  & \(+\) & \(-\) & \(-\)\\
III & \(-\) & \(-\) & \(+\)\\
IV  & \(-\) & \(+\) & \(-\)\\
\end{tabular}
\end{recordatoriobox}
\end{minipage}\hfill
\begin{minipage}[t]{0.50\textwidth}
\begin{recordatoriobox}{GRADOS Y RADIANES}
\[
    180^\circ=\pi\ \text{rad}
\]
\[
    \theta_{\mathrm{rad}}=\theta_{\mathrm{gra}}\frac{\pi}{180},
    \qquad
    \theta_{\mathrm{gra}}=\theta_{\mathrm{rad}}\frac{180}{\pi}.
\]
\end{recordatoriobox}
\end{minipage}

\begin{teoriabox}{FÓRMULAS QUE VAMOS A USAR}
\begin{multicols}{2}
\[
\boxed{\seno^2x+\cos^2x=1}
\]
\[
\boxed{1+\tan^2x=\sec^2x}
\]
\[
\boxed{1+\cot^2x=\csc^2x}
\]
\columnbreak
\[
\tan x=\frac{\seno x}{\cos x},
\qquad
\cot x=\frac{\cos x}{\seno x}
\]
\[
\sec x=\frac{1}{\cos x},
\qquad
\csc x=\frac{1}{\seno x}
\]
\end{multicols}
\vspace{-0.3cm}
\textbf{Para resolver una ecuación:} deja el seno, coseno o tangente solo; busca el valor en la tabla; marca todas las posiciones posibles en el círculo y conserva los ángulos que pide el ejercicio.
\end{teoriabox}

% ============================================================
% EJERCICIOS
% ============================================================
\section*{\color{azulFuerte}Entrenamiento progresivo}

\tituloSeccion{I. Orientarse en la vuelta}
\begin{enumerate}[label=\color{azulFuerte}\textbf{\arabic*.},leftmargin=*]
    \item Convierta a radianes:
    \[
        30^\circ,\qquad 135^\circ,\qquad 225^\circ,\qquad 330^\circ.
    \]
    \item Convierta a grados:
    \[
        \frac{2\pi}{3},\qquad \frac{5\pi}{4},\qquad \frac{11\pi}{6},\qquad -\frac{\pi}{3}.
    \]
    \item Encuentre el ángulo coterminal que pertenece a \([0,2\pi)\):
    \[
        \frac{13\pi}{6},\qquad -\frac{3\pi}{4},\qquad \frac{17\pi}{3},\qquad -\frac{5\pi}{2}.
    \]
    \item Indique el cuadrante y el ángulo de referencia:
    \[
        \frac{5\pi}{6},\qquad \frac{7\pi}{4},\qquad \frac{4\pi}{3},\qquad \frac{11\pi}{6}.
    \]
    \item Escriba las coordenadas \(P(\theta)=(\cos\theta,\seno\theta)\) para
    \[
        \theta=0,\qquad \frac{\pi}{2},\qquad \pi,\qquad \frac{3\pi}{2}.
    \]
    \item Indique solamente los signos de \(\seno\theta\), \(\cos\theta\) y \(\tan\theta\) si
    \[
        \theta=\frac{2\pi}{3},\qquad \frac{5\pi}{4},\qquad \frac{11\pi}{6}.
    \]
\end{enumerate}

\tituloSeccion{II. Leer valores exactos}
\begin{multicols}{2}
\begin{enumerate}[label=\color{azulFuerte}\textbf{\arabic*.},leftmargin=*,start=7]
    \item \(\seno\left(\frac{5\pi}{6}\right)\)
    \item \(\cos\left(\frac{7\pi}{4}\right)\)
    \item \(\tan\left(\frac{4\pi}{3}\right)\)
    \item \(\sec\left(\frac{2\pi}{3}\right)\)
    \item \(\csc\left(\frac{5\pi}{4}\right)\)
    \item Calcule:
    \[
    \seno\left(-\frac{\pi}{3}\right),\
    \cos\left(\frac{13\pi}{6}\right),\
    \tan\left(\frac{3\pi}{4}\right).
    \]
    \item Indique cuáles no existen:
    \[
    \tan\left(\frac{\pi}{2}\right),\
    \csc(\pi),\
    \sec\left(\frac{3\pi}{2}\right).
    \]
    \item Si \(P=(-12/13,5/13)\), indique el cuadrante y calcule las seis razones trigonométricas.
\end{enumerate}
\end{multicols}

\tituloSeccion{III. Usar identidades pitagóricas}
\begin{enumerate}[label=\color{azulFuerte}\textbf{\arabic*.},leftmargin=*,start=15]
    \item Si \(\seno\theta=-\frac35\) y \(\theta\) está en el IV cuadrante, calcule \(\cos\theta\) y \(\tan\theta\).
    \item Si \(\cos\theta=-\frac5{13}\) y \(\theta\) está en el III cuadrante, calcule \(\seno\theta\) y \(\tan\theta\).
    \item Si \(\tan\theta=-\frac8{15}\) y \(\theta\) está en el II cuadrante, calcule \(\seno\theta\) y \(\cos\theta\).
    \item Simplifique:
    \begin{multicols}{2}
    \begin{enumerate}[label=\alph*),leftmargin=*]
        \item \(\displaystyle\frac{1-\seno^2x}{\cos x}\)
        \item \(\displaystyle\frac{\sec^2x-1}{\tan x}\)
        \item \(\displaystyle(1-\cos^2x)\csc x\)
        \item \(\displaystyle\cos^2x(1+\tan^2x)\)
    \end{enumerate}
    \end{multicols}
    \item Demuestre que
    \[
        \tan x+\cot x=\sec x\,\csc x.
    \]
    \item Simplifique e indique las restricciones:
    \[
        \frac{1}{1+\seno x}+\frac{1}{1-\seno x}.
    \]
\end{enumerate}

\tituloSeccion{IV. Ecuaciones base}
\noindent En los ejercicios 21 al 30, resuelva en \(x\in[0,2\pi)\).

\begin{multicols}{2}
\begin{enumerate}[label=\color{azulFuerte}\textbf{\arabic*.},leftmargin=*,start=21,itemsep=0.65em]
    \item \(\seno x=\frac12\)
    \item \(\cos x=-\frac{\sqrt2}{2}\)
    \item \(\tan x=-\sqrt3\)
    \item
    \begin{enumerate}[label=\alph*),leftmargin=*]
        \item \(2\seno x+1=0\)
        \item \(2\cos x=\sqrt3\)
    \end{enumerate}
    \item \(4\seno^2x=3\)
    \item \(\tan^2x=1\)
    \item \(\seno x\cos x=0\)
    \item \(\seno^2x-\cos x=1\)
    \item \(2\seno^2x+3\cos x=3\)
    \item \(\sec^2x=4\)
\end{enumerate}
\end{multicols}

\tituloSeccion{V. Desafíos integradores}
\noindent En esta sección se combinan identidades, valores exactos y aplicaciones con ángulos notables.

\begin{desafiointegradorbox}{31. IDENTIDAD CON ÁNGULO AGUDO}
Sea \(\alpha\) un ángulo agudo tal que
\[
    \tan\alpha=2-\sec\alpha.
\]
Determine el valor exacto de \(\cos\alpha\).
\end{desafiointegradorbox}

\begin{desafiointegradorbox}{32. IDENTIDAD PITAGÓRICA Y ECUACIÓN}
Sea \(\alpha\) un ángulo agudo tal que
\[
    \frac{1}{\csc\alpha}=\frac{1+\cos\alpha}{2}.
\]
Determine el valor exacto de \(\seno\alpha\).
\end{desafiointegradorbox}

\begin{desafiointegradorbox}{33. SUMA DE ÁNGULOS}
Calcule el valor exacto de \(\seno(105^\circ)\) usando \(105^\circ=60^\circ+45^\circ\).
\end{desafiointegradorbox}

\begin{desafiointegradorbox}{34. CABLE TENSOR}
Una torre vertical de \(10\) m está instalada sobre una ladera. En el triángulo formado por la torre, el suelo y un cable tensor \(d\), el lado de \(10\) m está opuesto a \(30^\circ\) y el cable está opuesto a \(105^\circ\). Use la ley de senos y determine \(d\).
\end{desafiointegradorbox}

\begin{desafiointegradorbox}{35. AVIÓN Y DOS ÁNGULOS DE ELEVACIÓN}
Un avión se observa desde la base y desde la azotea de un edificio de \(20\) m. Los ángulos de elevación son \(60^\circ\) desde la base y \(45^\circ\) desde la azotea. Determine la altura del avión respecto del suelo.
\end{desafiointegradorbox}

\begin{desafiointegradorbox}{36. CAMBIO DE POSICIÓN}
Una persona observa la cima de una torre de \(50\) m con un ángulo de elevación de \(30^\circ\). Luego avanza en línea recta hacia la torre hasta observarla con un ángulo de \(60^\circ\). ¿Cuántos metros avanzó?
\end{desafiointegradorbox}

\begin{desafiointegradorbox}{37. ÁNGULO DOBLE DESDE UN PUNTO DEL CÍRCULO}
Si \(\alpha\) está en el II cuadrante y \(\cos\alpha=-\frac{12}{13}\), calcule
\[
    \seno(2\alpha)
    \qquad\text{y}\qquad
    \cos(2\alpha).
\]
\end{desafiointegradorbox}

\begin{desafiointegradorbox}{38. DESAFÍO DE CIERRE}
Resuelva en \(x\in[0,2\pi)\):
\[
    \frac{1}{1+\seno x}+\frac{1}{1-\seno x}=4.
\]
\end{desafiointegradorbox}

\begin{teoriabox}{AUTOEVALUACIÓN: RECONSTRUYE LA VUELTA}
\small
Sin mirar la tabla, rotule cada semirrecta con su ángulo en radianes. Luego escriba junto a cada punto sus coordenadas \((\cos\theta,\seno\theta)\).

\begin{center}
\begin{tikzpicture}[scale=3.15]
    \draw[->,gray!70] (-1.3,0)--(1.3,0) node[right] {\scriptsize \(x\)};
    \draw[->,gray!70] (0,-1.3)--(0,1.3) node[above] {\scriptsize \(y\)};
    \draw[azulFuerte,line width=1.1pt] (0,0) circle (1);
    \foreach \a in {0,30,45,60,90,120,135,150,180,210,225,240,270,300,315,330}{
        \draw[gray!45,thin] (0,0)--(\a:1);
        \filldraw[fill=white,draw=azulMedio,line width=0.6pt] (\a:1) circle (0.018);
    }
    \node[gray!70] at (0.55,0.72) {\scriptsize I};
    \node[gray!70] at (-0.55,0.72) {\scriptsize II};
    \node[gray!70] at (-0.55,-0.72) {\scriptsize III};
    \node[gray!70] at (0.55,-0.72) {\scriptsize IV};
\end{tikzpicture}
\end{center}

\textbf{Chequeo final:}\quad
\(\seno^2\theta+\cos^2\theta=\rule{1.2cm}{0.4pt}\),\qquad
\(P(\theta)=\bigl(\rule{1.5cm}{0.4pt},\rule{1.5cm}{0.4pt}\bigr)\).
\end{teoriabox}

% ============================================================
% SOLUCIONARIO
% ============================================================
\clearpage
\begin{center}
    \begin{tcolorbox}[
        colback=azulFuerte,
        colframe=azulFuerte,
        colupper=white,
        width=0.8\textwidth,
        center,
        arc=10pt]
        \centering\Large\textbf{SOLUCIONARIO}
    \end{tcolorbox}
\end{center}

\begin{solbox}{I. Orientarse en la vuelta}
\begin{enumerate}[label=\color{azulFuerte}\textbf{\arabic*.},leftmargin=*]
    \item \(\displaystyle \frac{\pi}{6},\ \frac{3\pi}{4},\ \frac{5\pi}{4},\ \frac{11\pi}{6}\).
    \item \(\displaystyle 120^\circ,\ 225^\circ,\ 330^\circ,\ -60^\circ\).
    \item \(\displaystyle \frac{\pi}{6},\ \frac{5\pi}{4},\ \frac{5\pi}{3},\ \frac{3\pi}{2}\).
    \item II y \(\pi/6\); IV y \(\pi/4\); III y \(\pi/3\); IV y \(\pi/6\).
    \item \((1,0),\ (0,1),\ (-1,0),\ (0,-1)\).
    \item II: \((+,-,-)\); III: \((-,-,+)\); IV: \((-,+,-)\), en el orden \((\seno,\cos,\tan)\).
\end{enumerate}
\end{solbox}

\begin{solbox}{II. Leer valores exactos}
\begin{multicols}{2}
\begin{enumerate}[label=\color{azulFuerte}\textbf{\arabic*.},leftmargin=*,start=7]
    \item \(\frac12\).
    \item \(\frac{\sqrt2}{2}\).
    \item \(\sqrt3\).
    \item \(-2\).
    \item \(-\sqrt2\).
    \item \(\displaystyle-\frac{\sqrt3}{2},\ \frac{\sqrt3}{2},\ -1\).
    \item Las tres expresiones no existen.
    \item Cuadrante II:
    \[
    \seno=\frac5{13},\quad
    \cos=-\frac{12}{13},\quad
    \tan=-\frac5{12},
    \]
    \[
    \csc=\frac{13}{5},\quad
    \sec=-\frac{13}{12},\quad
    \cot=-\frac{12}{5}.
    \]
\end{enumerate}
\end{multicols}
\end{solbox}

\begin{solbox}{III. Usar identidades pitagóricas}
\begin{enumerate}[label=\color{azulFuerte}\textbf{\arabic*.},leftmargin=*,start=15]
    \item
    \[
    \cos^2\theta=1-\frac9{25}=\frac{16}{25}.
    \]
    Como \(\theta\) está en IV, \(\cos\theta>0\). Así,
    \[
        \cos\theta=\frac45,\qquad \tan\theta=-\frac34.
    \]

    \item
    \[
    \seno^2\theta=1-\frac{25}{169}=\frac{144}{169}.
    \]
    En III, seno y coseno son negativos:
    \[
        \seno\theta=-\frac{12}{13},\qquad \tan\theta=\frac{12}{5}.
    \]

    \item Un triángulo de referencia con catetos \(8\) y \(15\) tiene hipotenusa \(17\). En II:
    \[
        \seno\theta=\frac8{17},\qquad \cos\theta=-\frac{15}{17}.
    \]

    \item
    \begin{multicols}{2}
    \begin{enumerate}[label=\alph*),leftmargin=*]
        \item \(\cos x\).
        \item \(\tan x\).
        \item \(\seno x\).
        \item \(1\).
    \end{enumerate}
    \end{multicols}

    \item
    \[
    \tan x+\cot x
    =\frac{\seno x}{\cos x}+\frac{\cos x}{\seno x}
    =\frac{\seno^2x+\cos^2x}{\seno x\cos x}
    =\frac1{\seno x\cos x}
    =\sec x\,\csc x.
    \]

    \item
    \[
    \frac{1}{1+\seno x}+\frac{1}{1-\seno x}
    =\frac{2}{1-\seno^2x}
    =\frac{2}{\cos^2x}
    =2\sec^2x.
    \]
    Restricciones: \(\seno x\neq\pm1\), equivalentes a \(\cos x\neq0\).
\end{enumerate}
\end{solbox}

\begin{solbox}{IV. Ecuaciones base}
\begin{enumerate}[label=\color{azulFuerte}\textbf{\arabic*.},leftmargin=*,start=21]
    \item \(\displaystyle x\in\left\{\frac{\pi}{6},\frac{5\pi}{6}\right\}\).
    \item \(\displaystyle x\in\left\{\frac{3\pi}{4},\frac{5\pi}{4}\right\}\).
    \item \(\displaystyle x\in\left\{\frac{2\pi}{3},\frac{5\pi}{3}\right\}\).
    \item
    \begin{enumerate}[label=\alph*),leftmargin=*]
        \item \(\displaystyle x\in\left\{\frac{7\pi}{6},\frac{11\pi}{6}\right\}\).
        \item \(\displaystyle x\in\left\{\frac{\pi}{6},\frac{11\pi}{6}\right\}\).
    \end{enumerate}
    \item \(\displaystyle x\in\left\{\frac{\pi}{3},\frac{2\pi}{3},\frac{4\pi}{3},\frac{5\pi}{3}\right\}\).
    \item \(\displaystyle x\in\left\{\frac{\pi}{4},\frac{3\pi}{4},\frac{5\pi}{4},\frac{7\pi}{4}\right\}\).
    \item \(\displaystyle x\in\left\{0,\frac{\pi}{2},\pi,\frac{3\pi}{2}\right\}\).
    \item Use \(\seno^2x=1-\cos^2x\):
    \[
        1-\cos^2x-\cos x=1
        \iff \cos x(\cos x+1)=0.
    \]
    Por tanto,
    \[
        x\in\left\{\frac{\pi}{2},\pi,\frac{3\pi}{2}\right\}.
    \]
    \item
    \[
    2(1-\cos^2x)+3\cos x=3
    \iff 2\cos^2x-3\cos x+1=0.
    \]
    Como \((2\cos x-1)(\cos x-1)=0\),
    \[
        x\in\left\{0,\frac{\pi}{3},\frac{5\pi}{3}\right\}.
    \]
    \item \(\sec^2x=4\iff\cos x=\pm\frac12\). Entonces,
    \[
        x\in\left\{\frac{\pi}{3},\frac{2\pi}{3},\frac{4\pi}{3},\frac{5\pi}{3}\right\}.
    \]
\end{enumerate}
\end{solbox}

\begin{solbox}{V. Desafíos integradores}
\begin{enumerate}[label=\color{azulFuerte}\textbf{\arabic*.},leftmargin=*,start=31,itemsep=1em]
    \item Como \(\alpha\) es agudo, \(\seno\alpha>0\) y \(\cos\alpha>0\). Reescribimos:
    \[
    \frac{\seno\alpha}{\cos\alpha}
    =2-\frac1{\cos\alpha}
    \iff
    \seno\alpha=2\cos\alpha-1.
    \]
    Al elevar al cuadrado y usar \(\seno^2\alpha=1-\cos^2\alpha\):
    \[
    1-\cos^2\alpha=(2\cos\alpha-1)^2
    \iff
    \cos\alpha(5\cos\alpha-4)=0.
    \]
    El valor \(\cos\alpha=0\) no corresponde a un ángulo agudo ni pertenece al dominio de la ecuación. Por tanto,
    \[
        \boxed{\cos\alpha=\frac45}.
    \]

    \item Como \(1/\csc\alpha=\seno\alpha\),
    \[
        2\seno\alpha=1+\cos\alpha
        \iff
        \cos\alpha=2\seno\alpha-1.
    \]
    Al cuadrar y usar \(\cos^2\alpha=1-\seno^2\alpha\):
    \[
        1-\seno^2\alpha=(2\seno\alpha-1)^2
        \iff
        \seno\alpha(5\seno\alpha-4)=0.
    \]
    Como \(\alpha\) es agudo,
    \[
        \boxed{\seno\alpha=\frac45}.
    \]

    \item
    \[
    \begin{aligned}
    \seno(105^\circ)
    &=\seno(60^\circ+45^\circ)\\
    &=\seno60^\circ\cos45^\circ+\cos60^\circ\seno45^\circ\\
    &=\frac{\sqrt3}{2}\frac{\sqrt2}{2}
      +\frac12\frac{\sqrt2}{2}
    =\boxed{\frac{\sqrt6+\sqrt2}{4}}.
    \end{aligned}
    \]

    \item Por la ley de senos,
    \[
        \frac{d}{\seno105^\circ}
        =\frac{10}{\seno30^\circ}.
    \]
    Luego,
    \[
        d=20\seno105^\circ
        =20\left(\frac{\sqrt6+\sqrt2}{4}\right)
        =\boxed{5(\sqrt6+\sqrt2)\ \text{m}}.
    \]

    \item Sea \(H\) la altura del avión y \(a\) la distancia horizontal al edificio. Entonces
    \[
        \tan60^\circ=\frac{H}{a},
        \qquad
        \tan45^\circ=\frac{H-20}{a}.
    \]
    De la segunda igualdad, \(a=H-20\). Sustituyendo:
    \[
        \sqrt3=\frac{H}{H-20}
        \iff
        H=30+10\sqrt3.
    \]
    La altura es
    \[
        \boxed{(30+10\sqrt3)\ \text{m}}.
    \]

    \item La distancia inicial a la torre es
    \[
        d_1=\frac{50}{\tan30^\circ}=50\sqrt3,
    \]
    y la distancia final es
    \[
        d_2=\frac{50}{\tan60^\circ}=\frac{50}{\sqrt3}.
    \]
    Por tanto, avanzó
    \[
        d_1-d_2
        =50\sqrt3-\frac{50}{\sqrt3}
        =\boxed{\frac{100\sqrt3}{3}\ \text{m}}.
    \]

    \item Como \(\alpha\) está en II, \(\seno\alpha=\frac5{13}\). Así,
    \[
        \seno(2\alpha)=2\seno\alpha\cos\alpha
        =2\left(\frac5{13}\right)\left(-\frac{12}{13}\right)
        =\boxed{-\frac{120}{169}},
    \]
    \[
        \cos(2\alpha)=\cos^2\alpha-\seno^2\alpha
        =\frac{144}{169}-\frac{25}{169}
        =\boxed{\frac{119}{169}}.
    \]

    \item La expresión está definida cuando \(\seno x\neq\pm1\). Usando el ejercicio 20:
    \[
        2\sec^2x=4
        \iff
        \cos^2x=\frac12.
    \]
    En \([0,2\pi)\),
    \[
        \boxed{x\in\left\{
        \frac{\pi}{4},\frac{3\pi}{4},
        \frac{5\pi}{4},\frac{7\pi}{4}
        \right\}}.
    \]
\end{enumerate}
\end{solbox}

\vspace{0.8cm}
\begin{center}
    \small\textbf{Usach Premium} - Nos vemos en la ayudantía.\\
    \textit{"Por y para estudiantes."}
\end{center}

\end{document}
